\documentclass[10pt,a4paper]{amsart}
\usepackage[margin=19mm]{geometry}
\usepackage{amsmath,amssymb,mathtools}
\usepackage[scr=rsfs,cal=euler]{mathalfa}
\usepackage{xfrac}
\usepackage[unicode,hidelinks,bookmarks=false]{hyperref}
\newcommand{\ie}{i.e.\ }
\newcommand{\cf}{cf.\ }
\newcommand{\resp}{respectively\ }
\newcommand{\Subsection}{\subsection}
\providecommand{\nobreakdash}{\nobreak\hbox{-}}
\setlength{\emergencystretch}{2em}
\multlinegap0pt
\usepackage{amssymb}
\usepackage{dsfont}
\usepackage{esint}
\newtheorem{lemma}{Lemma}
\title{Entropy-Minimizing Diffeomorphisms on a $G_2$-Manifold}
\author{Ollie Thakar}
\date{Source: arXiv:2602.07204v2 (CC BY 4.0)}
\begin{document}
\maketitle

\noindent\emph{Source and licence.} Entropy-Minimizing Diffeomorphisms on a $G_2$-Manifold. Version arXiv:2602.07204v2 (CC BY 4.0), distributed by arXiv under the Creative Commons Attribution 4.0 International licence, http://creativecommons.org/licenses/by/4.0/, as recorded in the arXiv licence metadata for this exact version and shown on the abstract page https://arxiv.org/abs/2602.07204v2. Full licence notice: \texttt{licenses/arxiv-2602.07204-CC-BY-4.0.txt}.\par\smallskip

\noindent\textit{Editorial scope.} Counted unit: the article's lemma lemma (source0.tex lines 226--228). The article's complete original proof of it is delivered with it (source0.tex lines 230--236, one \texttt{proof} environment of 7 lines). No mathematics is written, repaired or re-derived: the statement and the proof are the article's own lines, in the article's own notation, and no retained line is substituted or re-flowed.  No further source-local context is needed: every cross-reference of every retained line resolves inside the two delivered spans.  Nothing was omitted from any retained span, so the package carries no omission record.  No retained line contains a citation, so the wrapper carries no bibliography and no bibliography entry.  No retained line contains a figure environment, an image-inclusion command or drawing code, so no image is delivered and the package declares no asset dependency.  Editorial formatting only, in addition: an AMS \texttt{amsart} preamble that restates the article's own declarations and macros used by the retained lines, namely the environments \texttt{lemma} on separate counters, together with the standard packages those lines need.  The retained environments print in this document as Lemma 1 in order of appearance, whereas the article prints the same items with the numbering of its own full document.  The complete native source of the article as distributed in the arXiv e-print for this exact version is bundled unchanged as \texttt{sources/194157/source0.tex}.
\begin{lemma}
    The spectral radius of $(f_T)_*$ on $H_*(M)$ is $|\lambda|^2$ where $\lambda$ is the largest eigenvalue of $T.$ Moreover, it is achieved on $H_3(M).$
\end{lemma}

\begin{proof}
    $M$ is simply connected, so by Poincar\'e duality it is enough to consider $H_2(M)$ and $H_3(M).$ The second and third homology of $M$ is the direct sum of the elements of homology coming from $\mathbb{A}/\Gamma$ and from the resolutions. The action of $f_T$ on the homology permutes the factors coming from the resolutions, so we may consider the action of $f_T$ on $H_*(\mathbb{A}/\Gamma),$ which is the same as the induced action of the linear map $T$ on $H_*(\mathbb{A}/\Gamma).$

    We will compute using cohomology for ease: the cohomology of $\mathbb{A}/\Gamma$ is the $\Gamma$-invariant part of the cohomology of $\mathbb{A}.$ Hence, $H^2(\mathbb{A}/\Gamma)$ is generated by $\text{Re}(dz_i\wedge d\overline{z_j}),$ for $i\neq j,$ and $i,j\in\{1,2,3\}$. Likewise, $H^3(\mathbb{A}/\Gamma)$ is generated by $\text{Im}(dz_1\wedge dz_2\wedge dz_3)$, as well as $\text{Im}(dz_i\wedge d\overline{z_j}\wedge dx)$, for $i,j\in\{1,2,3\}$ potentially equal.

    Letting $\lambda_1,\lambda_2, \lambda_3$ be the three eigenvalues of $T$ in non-increasing order of their modulus, the spectral radius of $(f_T)_*$ is the maximum of $|\lambda_i\lambda_j|,$ potentially with $i=j,$ and $|\lambda_i\lambda_j\lambda_k|,$ where $i, j, k$ are all different. Since $T\in SL(3, \mathbb{Z}),$ we have that $\lambda_i\lambda_j\lambda_k=1$ and the maximum of $|\lambda_i\lambda_j|,$ potentially with $i=j,$ is given by $|\lambda_1|^2.$   
\end{proof}

\end{document}
